Narrow adversarial fine-tuning can broadly misalign a language model, yet the effect is famously inconsistent: many prompts, and many fine-tuning runs, show no drift at all. This inconsistency suggests that models hold \emph{latent robustness regions} that resist behavioral change. We ask whether this robustness is \emph{graded}: does per-prompt drift from the pre-fine-tuning baseline decrease with a prompt's semantic distance from the fine-tuning domain, forming low-drift ``basins'' on orthogonal tasks? We introduce a checkpoint-level \emph{trajectory-divergence} diagnostic. We LoRA-fine-tune a real instruct model (\smol) on the canonical \emph{Emergent Misalignment} \insecure-code dataset, checkpoint during training, and at each checkpoint greedily decode a graded 5-tier probe set spanning in-domain code to orthogonal math and science. Each probe's drift is a composite of embedding cosine distance and edit distance from its own pre-fine-tuning output, and we compare against a matched benign \secure-code control. We find no monotonic distance--drift relationship (Spearman $\rho=-0.08$, $p=0.78$); orthogonal-capability prompts drift the \emph{most} (0.55) and a code-safety Q\&A pocket drifts the \emph{least} (0.045), the opposite of the simple ``orthogonal $=$ robust'' prediction. Decisively, the adversarial and benign fine-tunes produce statistically indistinguishable drift (global mean 0.356 vs.\ 0.352), showing that at this scale trajectory divergence is dominated by generic fine-tuning style, not adversarial intent. We conclude that raw trajectory divergence is a valid robustness signal only when differenced against a matched benign-fine-tuning control, and that the one reproducible robustness pocket is content-type-specific (factual knowledge), not task distance.
